% 图 2-7 软件分层与调用关系（坐标单位 cm）
\begin{tikzpicture}[font=\footnotesize,
  L/.style={draw=softline, rounded corners=1.5pt, fill=maclav!70, align=center, inner sep=2.4pt,
            font=\scriptsize, minimum width=14.0cm, minimum height=1.25cm},
  b/.style={draw=softline, rounded corners=1.0pt, fill=white, align=center, inner sep=1.6pt,
            font=\scriptsize, text width=2.6cm, minimum height=1.0cm},
  ar/.style={-{Stealth[length=1.5mm]}, thick}]

\node[b, fill=macpeach!70] (v1) at (-5.2,0) {向量与期望值\\（参考实现、KAT）};
\node[b, fill=macpeach!70] (v2) at (-1.9,0) {仿真与回归\\（\texttt{run\_tests.py}）};
\node[b, fill=macpeach!70] (v3) at (1.4,0) {签核机检\\（\texttt{check\_io\_signoff}）};
\node[b, fill=macpeach!70] (v4) at (4.7,0) {PYNQ 端服务\\（\texttt{AsconHost}）};

\node[L] (host) at (0,-1.9) {\textbf{主机侧脚本与 PS 端服务}：向量生成、回归驱动、签核机检、Jupyter 交互};

\node[b, fill=macgreen!50] (f1) at (-5.2,-3.3) {\texttt{ascon\_rv32}\\纯软件};
\node[b, fill=macgreen!50] (f2) at (-1.9,-3.3) {\texttt{ascon\_hw\_rv32}\\MMIO 驱动};
\node[b, fill=macgreen!50] (f3) at (1.4,-3.3) {\texttt{ascon\_ise\_rv32}\\定制指令};
\node[b, fill=macamber] (f4) at (4.7,-3.3) {当前发布\\\texttt{pynq\_svc}};

\node[L] (fw) at (0,-5.2) {\textbf{PL 侧裸机固件}：前三套为历史基准；当前 \texttt{pynq\_svc} 含三模式，以 \texttt{mcycle} 计时};

\node[b, fill=macamber!70] (h1) at (-3.7,-6.6) {Ascon 置换核\\（SW / MMIO / ISE）};
\node[b, fill=macblue!55] (h2) at (0.4,-6.6) {AXI-Lite / UART\\字节 FIFO mailbox};
\node[b, fill=macgray] (h3) at (4.5,-6.6) {UART / CPI 计数\\（观测接口）};

\draw[ar] (v1) -- (host.north -| v1.south);
\draw[ar] (v2) -- (host.north -| v2.south);
\draw[ar] (v3) -- (host.north -| v3.south);
\draw[ar] (v4) -- (host.north -| v4.south);
\draw[ar] (f1) -- (fw.north -| f1.south);
\draw[ar] (f2) -- (fw.north -| f2.south);
\draw[ar] (f3) -- (fw.north -| f3.south);
\draw[ar] (f4) -- (fw.north -| f4.south);
\draw[ar] (host.south) -- node[right, font=\scriptsize] {AXI-Lite 帧协议} (fw.north);
\draw[ar] (h1) -- (fw.south -| h1.north);
\draw[ar] (h2) -- (fw.south -| h2.north);
\draw[ar] (h3) -- (fw.south -| h3.north);
\node[font=\scriptsize, anchor=west] at (5.9,-6.6) {硬件};
\node[font=\scriptsize, anchor=west] at (5.9,-5.2) {固件};
\node[font=\scriptsize, anchor=west] at (5.9,-1.9) {上位层};
\end{tikzpicture}
